% Curs 10 - Pagina 5 (6-XII-1994).
% Fragment pentru % Curs 10 - Pagina 5 (6-XII-1994).
% Fragment pentru % Curs 10 - Pagina 5 (6-XII-1994).
% Fragment pentru % Curs 10 - Pagina 5 (6-XII-1994).
% Fragment pentru \input{courses/pages/curs10-pagina5-ro}.
% Fara preambul si fara intrerupere fortata de pagina.

\subsection*{Curs 10 - Pagina 5}

\noindent\underline{Aplicație}: $G$ grup finit, $p\mid\operatorname{ord}(G)=n$.

\bigskip
\noindent\underline{Not.}
\[
A=\bigl\{(x_1,\ldots,x_p)\in
\underbrace{G\times\cdots\times G}_{p\text{ factori}}
\mid x_1x_2\cdots x_p=e\bigr\}.
\]

\noindent Atunci $|A|=n^{p-1}$:
\[
(x_1,\ldots,x_{p-1},
\bigl(\prod_{i=1}^{p-1}x_i\bigr)^{-1}).
\]

\medskip
\noindent Consider $\alpha=(1,2,\ldots,p)\in S_p$.

\medskip
\noindent Fie $H=\langle\alpha\rangle$, $\operatorname{ord}(\alpha)=p$.

\[
\Longrightarrow\quad H=\{e,\alpha,\ldots,\alpha^{p-1}\}.
\]
% In manuscris, litera din stanga multimii pare A, nu H;
% notatia H este necesara pentru subgrupul tocmai definit.

\noindent Dar $(x_1,\ldots,x_p)\in A$ implică:
\[
\left\{\begin{aligned}
&(x_2,x_3,\ldots,x_p,x_1)\in A,\\
&(x_3,x_4,\ldots,x_p,x_1,x_2)\in A,\\
&\quad\vdots\\
&(x_p,x_1,\ldots,x_{p-1})\in A.
\end{aligned}\right.
\]

\noindent $x_1x_2\cdots x_p=e$ și
\[
x_1x_2\cdots x_px_1=x_1
\quad\Longrightarrow\quad x_2\cdots x_px_1=e.
\]

\[
x_{\alpha(1)}x_{\alpha(2)}\cdots x_{\alpha(p)}=e,
\]
\[
x_{\alpha^{p-1}(1)}x_{\alpha^{p-1}(2)}
\cdots x_{\alpha^{p-1}(p)}=e.
\]

\bigskip
\noindent Fie $H=\langle\alpha\rangle$. Aplic lema 1.

\[
\Longrightarrow\quad |A|\equiv |A_0|\pmod p.
\]

\noindent Cum $|A|=n^{p-1}$ și $p\mid n$,
\[
p\mid |A|\quad\Longrightarrow\quad p\mid |A_0|.
\]

\noindent $|A_0|\geq 1$.

\medskip
\noindent\underline{Ex.}:
\[
(x_1,\ldots,x_p)\in A_0
\quad\Longleftrightarrow\quad
x_1=\cdots=x_p=a
\quad\Longrightarrow\quad a^p=e.
\]

\[
\Longrightarrow\quad
A_0=\bigl\{(\underbrace{a,\ldots,a}_{p\text{ ori}})
\mid a\in G,\ a^p=e\bigr\}.
\]

\[
\Longrightarrow\quad |A_0|\geq p
\quad\Longrightarrow\quad
\text{există cel puțin }p\text{ elemente cu }a^p=e.
\]

\[
\Longrightarrow\quad (\exists)\,a\neq 1
\ \text{a.î.}\ a^p=e
\qquad (\operatorname{ord}(a)=p).
\]
.
% Fara preambul si fara intrerupere fortata de pagina.

\subsection*{Curs 10 - Pagina 5}

\noindent\underline{Aplicație}: $G$ grup finit, $p\mid\operatorname{ord}(G)=n$.

\bigskip
\noindent\underline{Not.}
\[
A=\bigl\{(x_1,\ldots,x_p)\in
\underbrace{G\times\cdots\times G}_{p\text{ factori}}
\mid x_1x_2\cdots x_p=e\bigr\}.
\]

\noindent Atunci $|A|=n^{p-1}$:
\[
(x_1,\ldots,x_{p-1},
\bigl(\prod_{i=1}^{p-1}x_i\bigr)^{-1}).
\]

\medskip
\noindent Consider $\alpha=(1,2,\ldots,p)\in S_p$.

\medskip
\noindent Fie $H=\langle\alpha\rangle$, $\operatorname{ord}(\alpha)=p$.

\[
\Longrightarrow\quad H=\{e,\alpha,\ldots,\alpha^{p-1}\}.
\]
% In manuscris, litera din stanga multimii pare A, nu H;
% notatia H este necesara pentru subgrupul tocmai definit.

\noindent Dar $(x_1,\ldots,x_p)\in A$ implică:
\[
\left\{\begin{aligned}
&(x_2,x_3,\ldots,x_p,x_1)\in A,\\
&(x_3,x_4,\ldots,x_p,x_1,x_2)\in A,\\
&\quad\vdots\\
&(x_p,x_1,\ldots,x_{p-1})\in A.
\end{aligned}\right.
\]

\noindent $x_1x_2\cdots x_p=e$ și
\[
x_1x_2\cdots x_px_1=x_1
\quad\Longrightarrow\quad x_2\cdots x_px_1=e.
\]

\[
x_{\alpha(1)}x_{\alpha(2)}\cdots x_{\alpha(p)}=e,
\]
\[
x_{\alpha^{p-1}(1)}x_{\alpha^{p-1}(2)}
\cdots x_{\alpha^{p-1}(p)}=e.
\]

\bigskip
\noindent Fie $H=\langle\alpha\rangle$. Aplic lema 1.

\[
\Longrightarrow\quad |A|\equiv |A_0|\pmod p.
\]

\noindent Cum $|A|=n^{p-1}$ și $p\mid n$,
\[
p\mid |A|\quad\Longrightarrow\quad p\mid |A_0|.
\]

\noindent $|A_0|\geq 1$.

\medskip
\noindent\underline{Ex.}:
\[
(x_1,\ldots,x_p)\in A_0
\quad\Longleftrightarrow\quad
x_1=\cdots=x_p=a
\quad\Longrightarrow\quad a^p=e.
\]

\[
\Longrightarrow\quad
A_0=\bigl\{(\underbrace{a,\ldots,a}_{p\text{ ori}})
\mid a\in G,\ a^p=e\bigr\}.
\]

\[
\Longrightarrow\quad |A_0|\geq p
\quad\Longrightarrow\quad
\text{există cel puțin }p\text{ elemente cu }a^p=e.
\]

\[
\Longrightarrow\quad (\exists)\,a\neq 1
\ \text{a.î.}\ a^p=e
\qquad (\operatorname{ord}(a)=p).
\]
.
% Fara preambul si fara intrerupere fortata de pagina.

\subsection*{Curs 10 - Pagina 5}

\noindent\underline{Aplicație}: $G$ grup finit, $p\mid\operatorname{ord}(G)=n$.

\bigskip
\noindent\underline{Not.}
\[
A=\bigl\{(x_1,\ldots,x_p)\in
\underbrace{G\times\cdots\times G}_{p\text{ factori}}
\mid x_1x_2\cdots x_p=e\bigr\}.
\]

\noindent Atunci $|A|=n^{p-1}$:
\[
(x_1,\ldots,x_{p-1},
\bigl(\prod_{i=1}^{p-1}x_i\bigr)^{-1}).
\]

\medskip
\noindent Consider $\alpha=(1,2,\ldots,p)\in S_p$.

\medskip
\noindent Fie $H=\langle\alpha\rangle$, $\operatorname{ord}(\alpha)=p$.

\[
\Longrightarrow\quad H=\{e,\alpha,\ldots,\alpha^{p-1}\}.
\]
% In manuscris, litera din stanga multimii pare A, nu H;
% notatia H este necesara pentru subgrupul tocmai definit.

\noindent Dar $(x_1,\ldots,x_p)\in A$ implică:
\[
\left\{\begin{aligned}
&(x_2,x_3,\ldots,x_p,x_1)\in A,\\
&(x_3,x_4,\ldots,x_p,x_1,x_2)\in A,\\
&\quad\vdots\\
&(x_p,x_1,\ldots,x_{p-1})\in A.
\end{aligned}\right.
\]

\noindent $x_1x_2\cdots x_p=e$ și
\[
x_1x_2\cdots x_px_1=x_1
\quad\Longrightarrow\quad x_2\cdots x_px_1=e.
\]

\[
x_{\alpha(1)}x_{\alpha(2)}\cdots x_{\alpha(p)}=e,
\]
\[
x_{\alpha^{p-1}(1)}x_{\alpha^{p-1}(2)}
\cdots x_{\alpha^{p-1}(p)}=e.
\]

\bigskip
\noindent Fie $H=\langle\alpha\rangle$. Aplic lema 1.

\[
\Longrightarrow\quad |A|\equiv |A_0|\pmod p.
\]

\noindent Cum $|A|=n^{p-1}$ și $p\mid n$,
\[
p\mid |A|\quad\Longrightarrow\quad p\mid |A_0|.
\]

\noindent $|A_0|\geq 1$.

\medskip
\noindent\underline{Ex.}:
\[
(x_1,\ldots,x_p)\in A_0
\quad\Longleftrightarrow\quad
x_1=\cdots=x_p=a
\quad\Longrightarrow\quad a^p=e.
\]

\[
\Longrightarrow\quad
A_0=\bigl\{(\underbrace{a,\ldots,a}_{p\text{ ori}})
\mid a\in G,\ a^p=e\bigr\}.
\]

\[
\Longrightarrow\quad |A_0|\geq p
\quad\Longrightarrow\quad
\text{există cel puțin }p\text{ elemente cu }a^p=e.
\]

\[
\Longrightarrow\quad (\exists)\,a\neq 1
\ \text{a.î.}\ a^p=e
\qquad (\operatorname{ord}(a)=p).
\]
.
% Fara preambul si fara intrerupere fortata de pagina.

\subsection*{Curs 10 - Pagina 5}

\noindent\underline{Aplicație}: $G$ grup finit, $p\mid\operatorname{ord}(G)=n$.

\bigskip
\noindent\underline{Not.}
\[
A=\bigl\{(x_1,\ldots,x_p)\in
\underbrace{G\times\cdots\times G}_{p\text{ factori}}
\mid x_1x_2\cdots x_p=e\bigr\}.
\]

\noindent Atunci $|A|=n^{p-1}$:
\[
(x_1,\ldots,x_{p-1},
\bigl(\prod_{i=1}^{p-1}x_i\bigr)^{-1}).
\]

\medskip
\noindent Consider $\alpha=(1,2,\ldots,p)\in S_p$.

\medskip
\noindent Fie $H=\langle\alpha\rangle$, $\operatorname{ord}(\alpha)=p$.

\[
\Longrightarrow\quad H=\{e,\alpha,\ldots,\alpha^{p-1}\}.
\]
% In manuscris, litera din stanga multimii pare A, nu H;
% notatia H este necesara pentru subgrupul tocmai definit.

\noindent Dar $(x_1,\ldots,x_p)\in A$ implică:
\[
\left\{\begin{aligned}
&(x_2,x_3,\ldots,x_p,x_1)\in A,\\
&(x_3,x_4,\ldots,x_p,x_1,x_2)\in A,\\
&\quad\vdots\\
&(x_p,x_1,\ldots,x_{p-1})\in A.
\end{aligned}\right.
\]

\noindent $x_1x_2\cdots x_p=e$ și
\[
x_1x_2\cdots x_px_1=x_1
\quad\Longrightarrow\quad x_2\cdots x_px_1=e.
\]

\[
x_{\alpha(1)}x_{\alpha(2)}\cdots x_{\alpha(p)}=e,
\]
\[
x_{\alpha^{p-1}(1)}x_{\alpha^{p-1}(2)}
\cdots x_{\alpha^{p-1}(p)}=e.
\]

\bigskip
\noindent Fie $H=\langle\alpha\rangle$. Aplic lema 1.

\[
\Longrightarrow\quad |A|\equiv |A_0|\pmod p.
\]

\noindent Cum $|A|=n^{p-1}$ și $p\mid n$,
\[
p\mid |A|\quad\Longrightarrow\quad p\mid |A_0|.
\]

\noindent $|A_0|\geq 1$.

\medskip
\noindent\underline{Ex.}:
\[
(x_1,\ldots,x_p)\in A_0
\quad\Longleftrightarrow\quad
x_1=\cdots=x_p=a
\quad\Longrightarrow\quad a^p=e.
\]

\[
\Longrightarrow\quad
A_0=\bigl\{(\underbrace{a,\ldots,a}_{p\text{ ori}})
\mid a\in G,\ a^p=e\bigr\}.
\]

\[
\Longrightarrow\quad |A_0|\geq p
\quad\Longrightarrow\quad
\text{există cel puțin }p\text{ elemente cu }a^p=e.
\]

\[
\Longrightarrow\quad (\exists)\,a\neq 1
\ \text{a.î.}\ a^p=e
\qquad (\operatorname{ord}(a)=p).
\]
